\documentclass[11pt]{article}
\usepackage{amsmath,amssymb,amsthm,mathtools,booktabs}
\usepackage[margin=1in]{geometry}
\usepackage[colorlinks=true,linkcolor=blue,urlcolor=blue,citecolor=blue]{hyperref}
\usepackage{xurl}
\newcommand{\Q}{\mathbb Q}
\newcommand{\Z}{\mathbb Z}
\newcommand{\F}{\mathbb F}
\newcommand{\vp}{v_p}
\newtheorem{theorem}{Theorem}
\newtheorem{lemma}[theorem]{Lemma}
\newtheorem{corollary}[theorem]{Corollary}
\theoremstyle{remark}\newtheorem{remark}[theorem]{Remark}
\title{A one-parameter pole criterion and an obstruction\ to shifted abscissa menus}
\author{Project research note}
\date{20 November 2026, revision v82}
\begin{document}
\maketitle
\begin{abstract}
A genus-one equation with one free pole parameter gives the local
inequality required by the existing eleven-variable construction. Every
prepared scalar fibre is everywhere locally soluble. We give an exact
nonsquare-\(u\) counterexample to rational solubility, even after allowing
all three shifted elliptic abscissas and the standard sixteen scalar
multipliers. The rational-rank bound is certified by algebraic descent.
The counterexample is prepared parameter data, not a claimed value of the
actual contact detector. No smaller definition of \(\Z\) is established.
\end{abstract}

This note develops the source notes of \textbf{Gabe Conant with Astra 6.2}.
It separates a proved soundness criterion from its open actual-input
completeness requirement. All references to the eleven-variable source
mean the guarded construction in \cite{eleven}; its ordinary rational
\(abc\) premise is retained.

\section{A sound one-parameter criterion}
Let \(u,V\in\Q\), with \(u(1+u)V\ne0\), and put
\begin{equation}\label{eq:data}
 w=\frac{u}{(1+u)^2},\qquad d=7w,\qquad
 H=V^2+d,\qquad N=dH.
\end{equation}
The proposed replacement for the two intrinsic pole parameters is
\begin{equation}\label{eq:cover}
 C_{d,V}:\qquad Y^2=d(s^4-H^2).
\end{equation}
Only \(s\) would be a free rational coordinate; \(Y\) is a terminal
quadratic root in the finite algebra. When a nonzero point is needed for
finite avoidance, we impose \(sY\ne0\).

\begin{lemma}\label{lem:sound}
Let \(p\notin\{2,7\}\) and suppose \(\mu=\vp(u)>0\) is odd.
A rational solution of \eqref{eq:cover} implies
\begin{equation}\label{eq:needed}
 2\vp(V)<\mu.
\end{equation}
In particular \(\vp(4d/V^2)>0\) is odd.
\end{lemma}
\begin{proof}
Here \(\vp(d)=\mu\). Equality \(2\vp(V)=\mu\) is impossible.
If \(2\vp(V)>\mu\), then \(\vp(H)=\mu\). For \(s\ne0\), the two
terms on the right of \eqref{eq:cover} have valuations
\(\mu+4\vp(s)\) and \(3\mu\). These are distinct odd integers, since
equality would require \(\vp(s)=\mu/2\). Their sum is not a square.
For \(s=0\), the single nonzero term also has odd valuation. This includes
\(Y=0\), since cancellation of the two terms was excluded.
\end{proof}

The displayed inequality is the scalar conclusion used by the current
source to extract its six local tests. The derivative contact radical and
the subsequent height comparison are unchanged. This gives the following
precise reduction criterion, rather than an additional consequence of
ordinary \(abc\).

\begin{theorem}[Actual-input criterion]\label{thm:criterion}
Retain ordinary rational \(abc\) and all premises and base guards of
\cite{eleven}. Suppose that for every nonzero integer input there is a
choice of its actual prepared source witnesses such that
\eqref{eq:cover} has a rational point with \(sY\ne0\). Then replacing the
intrinsic pole by \eqref{eq:cover}, and using the same finite-algebra
compression, gives a defining polynomial with ten rational coordinates.
The additional actual-input rational-point hypothesis is not proved here.
\end{theorem}
\begin{proof}
Lemma~\ref{lem:sound} gives precisely the existing local extraction at each
positive odd contact valuation; the existing height argument therefore
proves soundness. The polynomial change
\begin{equation}\label{eq:map}
 x=d s^2,\qquad y=d sY
\end{equation}
maps \eqref{eq:cover} to \(E_N:y^2=x^3-N^2x\).
The hypotheses \(sY\ne0\) give a point outside the rational two-torsion.
A congruent-number curve has rational torsion \(\Z/2\Z\times\Z/2\Z\),
so this point has infinite order. Its odd multiples retain the same
abscissa square class, and hence give infinitely many rational \(s\)
on the cover. One can therefore avoid the finitely many terminal zeros
at all point-algebra characters. Each such polynomial in \(s\) is nonzero:
its leading coefficient is the same nonzero scalar \(d\).
The already guarded point algebra has rank \(64\); the terminal root
makes rank \(128\), and one primitive-element coordinate suffices as in
\cite{encoder}. The count is
\[
 \underbrace{1}_{\text{input}}+
 \underbrace{6}_{\text{three point blocks}}+
 \underbrace{1}_{\text{contact slope}}+
 \underbrace{1}_{s}+
 \underbrace{1}_{\text{encoder}}=10.
\]
Zero input can be added by the same finite exceptional-fibre convention
as in the source. This proof uses the stated rational-point hypothesis;
local solubility does not supply it.
\end{proof}

\begin{corollary}[One quadratic progression is sufficient]\label{cor:progression}
Fix one admissible source constant \(B\) for which the preceding
soundness argument holds. It is enough to prove the additional
actual-input hypothesis of Theorem~\ref{thm:criterion} at source inputs
\(cm^2+b\) for all but finitely many \(m\in\Z\), where
\(c>0\) is one fixed squarefree integer and \(b\in\Z\) is fixed.
The source constant \(B\) must be the same throughout this progression.
For example, actual witnesses at \(2m^2+1\) would suffice; their
existence is not asserted here.
\end{corollary}
\begin{proof}
Independently of completeness, the guarded construction in the preceding
proof gives a polynomial \(P_B(T,\mathbf Y)\) in ten coordinates whose
rational projection is contained in \(\Z\). Substitute
\(T=cX^2+b\). If \(X=a/q\) is reduced, with \(q>0\), and this
substituted polynomial has a rational zero, then \(ca^2/q^2\in\Z\).
Thus \(q^2\mid c\), and squarefreeness gives \(q=1\).
The projection remains sound for \(\Z\).

Let \(F\subset\Z\) contain the finitely many omitted integers, including
any integer whose source input \(cm^2+b\) is zero. Then
\[
 \left(\prod_{m\in F}(X-m)\right)P_B(cX^2+b,\mathbf Y)
\]
defines all of \(\Z\) and has the same ten coordinates. This is the
elementary rational-input version of the power-shift and finite-exception
method in \cite{powers}. It weakens the open global realization target;
it does not establish that target on any progression.
\end{proof}

\section{Everywhere-local solubility}
\begin{theorem}\label{thm:local}
Suppose \(d>0\), \(V\ne0\), \(d\in7\Q_2^{\times2}\), and
\(e=4d/V^2\) has square reduced denominator. Then \(C_{d,V}\) has a
point over every completion of \(\Q\). At every completion the points
can be chosen with \(sY\ne0\).
\end{theorem}
\begin{proof}
Write \(H=V^2+d>0\). At an odd prime \(p\), let \(\mu=\vp(d)\).
If \(\mu\) is odd, including a negative odd value, the denominator
condition on \(e\) implies \(\mu-2\vp(V)>0\). Consequently
\(H/V^2\in1+p\Z_p\) is a square. A point with \(s^2=H\) and \(Y=0\)
is smooth; its local neighbourhood contains points with \(sY\ne0\).

Now let \(\mu=2m\). If \(d\) is a local square, a sufficiently negative value
of \(\vp(s)\) makes \(d s^4(1-H^2/s^4)\) a square.
If \(-d\) is a local square, \(s=0\) supplies a smooth point, which can
again be perturbed. Suppose neither is a square. If \(\vp(V)<m\),
then \(H\) is a square as before. Otherwise \(\vp(H)=2m\): higher
valuation would force the residue of \(-d p^{-2m}\) to be a square.
Write \(d=p^{2m}D\), \(H=p^{2m}H_0\), and set
\(s=p^m s_0\), \(Y=p^{3m}Y_0\). This gives
\[
 Y_0^2=D(s_0^4-H_0^2),\qquad D,H_0\in\Z_p^\times.
\]
Its smooth projective reduction is a genus-one curve: the four roots
of the quartic are distinct over \(\overline\F_p\). The genus-one Weil
bound gives an \(\F_p\)-point. Its points at infinity are not rational
over \(\F_p\), since \(D\) is a nonsquare, so an affine smooth point
lifts. A local disk avoids the finitely many zeros of \(sY\).

At \(2\), \(-d\in\Q_2^{\times2}\), so \(s=0\) supplies a smooth
point and the same perturbation applies. Over \(\mathbb R\), choose
\(s>\sqrt H\). No assumption on the denominator of \(d\) itself was
used.
\end{proof}

\section{The shifted-abscissa menu}
In addition to \eqref{eq:cover}, setting \(x+N=d s^2\) or
\(x-N=d s^2\) gives the two covers
\begin{align}
 Y^2&=d(s^2-H)(s^2-2H),\label{eq:plus}\\
 Y^2&=d(s^2+H)(s^2+2H).\label{eq:minus}
\end{align}
The maps have \(y=d sY\) and respectively
\(x=d s^2-N\), \(x=d s^2+N\).
These covers, or finite menus of fixed square-class multipliers, cost
no additional free coordinate.

\begin{lemma}\label{lem:shifts}
At an odd prime, if \(\vp(N)\) is even, every nonzero one of
\(x,x-N,x+N\) at a rational point of \(E_N\) has even valuation.
Thus requiring any of these coordinates to lie in
\(\kappa d\Q^{\times2}\), with \(\vp(\kappa)\) even, retains the
obstruction of Lemma~\ref{lem:sound} when \(\vp(d)\) is odd and \(2\vp(V)>\vp(d)\).
\end{lemma}
\begin{proof}
Write \(N=p^{2j}N_0\) with \(N_0\) a unit and scale
\(x=p^{2j}X\), \(y=p^{3j}Y\). If \(\vp(X)<0\), the equation
\(Y^2=X(X-N_0)(X+N_0)\) forces \(3\vp(X)\) even.
If \(X\) is integral, at most one of the three factors has positive
valuation; its valuation is even. The nonzero coordinates at a
 two-torsion point are \(\pm N_0,\pm2N_0\), which are units.
At a bad candidate valuation, \(\vp(H)=\vp(d)\), so
\(\vp(N)=2\vp(d)\), whereas \(\kappa d\) has odd valuation.
\end{proof}

\section{An exact obstruction with nonsquare input}
\begin{theorem}\label{thm:counter}
There is prepared rational data \eqref{eq:data} with nonsquare \(u\),
\(0<V<2\), \(0<e=4d/V^2<1\), square reduced denominators for both
\(d,e\), and
\[
 d,e\in7\Q_2^{\times2},\qquad
 v_2(u)=-8,\quad v_2(d)=8,\quad v_2(V)=1,\quad v_2(e)=8.
\]
Nevertheless no rational point on \(E_{dH}\) has any nonzero coordinate
\(x,x-dH,x+dH\) in \(\kappa d\Q^{\times2}\) for any nonzero rational
\(\kappa\) supported on \(\{2,3,5,7\}\). The exclusion remains true
after multiplying \(\kappa\) by any integral power of \(dH\).
\end{theorem}
\begin{proof}
Set
\[
 \beta=119=7\cdot17,\qquad \eta=52985=5\cdot10597,
 \qquad N_0=\beta\eta=6305215.
\]
Both \(17\) and \(10597\) are prime. The exact algebraic descent
computation in the accompanying checker gives
\begin{equation}\label{eq:rank}
 \operatorname{rank}E_{N_0}(\Q)=1,
 \qquad E_{N_0}(\Q)_{\mathrm{tors}}\cong(\Z/2\Z)^2.
\end{equation}
In PARI/GP notation the complete rank output is
\begin{verbatim}
[1,1,2,[[-6002500,4729026750]]].
\end{verbatim}
The first two entries are unconditional lower and upper rank bounds
from two-descent and the Cassels pairing \cite{pari}. This is not an
analytic-rank approximation. The displayed point \(R\) satisfies
\begin{equation}\label{eq:point}
 x(R)=-2450^2,\quad x(R)-N_0=-35\cdot593^2,\quad
 x(R)+N_0=35\cdot93^2.
\end{equation}

Use the Kummer map \(\delta(P)=(x(P),x(P)-N_0)\) modulo squares,
whose kernel is \(2E_{N_0}(\Q)\); see \cite[Section~3.2]{burungale}.
The two torsion generators have classes \((-1,-N_0)\) and \((N_0,2)\).
The class of \(R\) is \((-1,-35)\), and its normalized class is
\((1,180149)\), distinct from the normalized torsion class. It follows
from \eqref{eq:rank} that these three classes generate the entire
Kummer image. All nonzero coordinates \(x,x-N_0,x+N_0\), including
the exceptional coordinates at two-torsion, therefore belong to
\begin{equation}\label{eq:G}
 \mathcal G=\langle-1,2,35,N_0\rangle
       \subset\Q^\times/\Q^{\times2}.
\end{equation}
Every class in \(\mathcal G\) has equal valuation parity at \(17\) and
\(10597\). The class of \(\kappa\beta\) has unequal parity there for
every \(\kappa\) in the theorem. Multiplication by \(N_0\) changes
both parities and cannot fix this discrepancy.

For explicit prepared data, use
\begin{gather*}
 213^2+119\cdot8^2=52985,\qquad
 t=\frac{19}{16},\quad u=17t^2=\frac{6137}{256},\quad
 r=\frac{t}{1+u}=\frac{304}{6393},\\
 d=119r^2=\frac{10997504}{40870449},\qquad
 V=\frac{213}{8}r=\frac{2698}{2131},\qquad
 H=52985\left(\frac r8\right)^2.
\end{gather*}
Then \(d=7u/(1+u)^2\), \(H=V^2+d\), and
\[
 dH=N_0\left(\frac{r^2}{8}\right)^2,
 \qquad e=\frac{4d}{V^2}=\frac{30464}{45369}
                    =119\left(\frac{16}{213}\right)^2.
\]
All bounds, denominator assertions and valuations in the statement follow
by exact rational arithmetic. The scaling from \(E_{dH}\) to \(E_{N_0}\)
multiplies every shifted abscissa by a rational square. Its target class
is \(d\Q^{\times2}=119\Q^{\times2}\), excluded by \eqref{eq:G}.
\end{proof}

The selected-prime residue pattern is already present in this example:
\(17\equiv1\pmod8\), \((7/17)=-1\), and \(v_{17}(u)=1\).
Also \(\eta\equiv1\pmod8\), and \(\eta\) is a square in both
\(\Q_7\) and \(\Q_{17}\). Thus even the idealized single-leaf star
pattern does not imply rational realization. The exclusion contains the
usual sixteen multipliers
\(\{\pm1,\pm2,\pm3,\pm5,\pm6,\pm7,\pm10,\pm14\}\).

\begin{remark}[Depth and real size]
The failure is not confined to the displayed two-adic depth or real size.
The conic \(a^2+119b^2=52985\) is rational and has the point \((213,8)\).
Weak approximation on this conic lets \(q=a/b\) be simultaneously positive
and arbitrarily large over \(\mathbb R\), and have \(v_2(q)=-k\) for
any prescribed \(k\ge2\). Indeed its two-adic branch near
\((\sqrt{52985},0)\) has precisely these points. Then
\(e=4\cdot119/q^2\) is arbitrarily small and has
\(v_2(e)=2+2k\). Its reduced denominator remains square: at a prime
of \(119\), the conic equation and coprimality with \(52985\) force
\(a\) integral and a unit, \(b\) integral, hence \(v_p(q)\le0\).
At other odd primes the valuations of \(e\) are even.

Now choose a rational \(t\) large and positive over \(\mathbb R\), with
\(v_2(t)=-(k+1)\), and a unit at \(17\). With
\(u=17t^2\), \(r=t/(1+u)\), \(d=119r^2\), and \(V=qr\), the same
curve square class \(N_0\) and the same global obstruction persist.
One obtains \(0<V<2\), \(v_2(V)=1\), and \(v_{17}(u)=1\).
Choosing \(v_7(t)\ge1\) as well makes the reduced denominator
of \(d\) square. This variation does not assert actual contact realization.
\end{remark}

\section{Verification and the remaining question}
The package \texttt{check\_genus\_one.py} reconstructs the three curve maps,
checks the prepared tuple and norm identity, computes both algebraic rank
bounds, enumerates all eight Kummer classes, and verifies the excluded
square classes. It independently applies \texttt{ellrank} to the
\(2\)-isogenous curve \(y^2=x^3+4N_0^2x\), obtaining the same rank
interval \([1,1]\); this is a second specialization of the same library
algorithm. The default check is read-only against \texttt{certificate.json}.
The recorded runtime is PARI/GP~2.17.4 through PassageMath~10.8.14.

The local theorem is proved above, not inferred from checking finitely
many places. The counterexample is not claimed to equal the derivative
contact value and trace supplied by \cite{eleven}. The new ten-coordinate
criterion therefore remains open on those actual inputs. Additional
local preparation, or the standard shifted-abscissa menu alone, does not
supply its missing global rational-point theorem.

\begin{thebibliography}{9}
\bibitem{eleven}
Gabe Conant with Astra 6.2, source notes and project audit,
\emph{Twelve contact shifts at the first mixed-power scale}, v78,
\nolinkurl{papers/conditional_few_variables/twelve_contact_first_scale_20261119_v78.tex}.

\bibitem{encoder}
Project v81 finite-algebra audit,
\emph{A ten-coordinate definition of the rational subfield by a guarded finite algebra},
\nolinkurl{papers/number_fields/ten_coordinate_rationality_20261119_v81.tex}.
The field-rationality theorem itself is not being used to define
\(\Z\subset\Q\); only its algebraic compression is used.

\bibitem{powers}
Project finite-exception audit,
\emph{Power shifts and finite exceptions}, v69,
\nolinkurl{papers/number_fields/power_shifts_and_finite_exceptions_20261117_v69.tex}.

\bibitem{burungale}
A.~Burungale, \emph{A Diophantine definition of \(\Z\) in \(\Q\)},
arXiv:2611.16130v1, Section~3.2,
\url{https://arxiv.org/html/2611.16130v1}.

\bibitem{pari}
The PARI Group, \emph{PARI/GP function catalogue: Elliptic curves},
\texttt{ellrank}, documentation of unconditional rank bounds,
\url{https://pari.math.u-bordeaux.fr/dochtml/html/Elliptic_curves.html#ellrank}.
\end{thebibliography}
\end{document}
